\scanpage{32}
D'autre part, si $F$ est un faisceau constant, de la forme $I_X$, \deleted{\unsure{on}} alors
\diagram{\begin{tikzcd}[column sep=huge]
\operatorname{Fr}_X^*(I_X)\simeq I_X \arrow[r,leftrightarrow,"\operatorname{Fr}_{F/X}^*=\mathrm{id}\,I_X"] \arrow[r,phantom,"\unsure{\rightarrow}" description] & I_X
\end{tikzcd}}
2) \deleted{Ces conditions} \deleted{\unsure{Pour $X$ fixé}} Ces conditions (pour $X,F$ variables) et la fonctorialité en $F$, \uline{caractérisent} les hom $\operatorname{Fr}_{F/X}$.
\margin{N.B. la condition relative à $I_X$ est \unsure{même} inutile, car \unsure{on se ramène au} cas où $I$ -- un seul élément, et alors c'est trivial, car le faisceau final \unsure{a} un seul endomorphisme !}
\inserted{il suffit de prendre $X'$ étale sur $X$ --}

Prouvons l'assertion $2^\circ)$. Comme $\forall F$ \deleted{\unsure{étale}} est quotient d'un $G$ représentable, \deleted{\unsure{Alors}} on est ramené à déterminer $\operatorname{Fr}_{F/X}^*$ pour $F$ \deleted{\unsure{étale}} représentable.
\margin{\diagram{\begin{tikzcd}[column sep=large,row sep=small]\operatorname{Fr}_X^*(G) \arrow[d] & G \arrow[l] \arrow[d,"\text{épim}"]\\\operatorname{Fr}_X^*(F) & F \arrow[l]\end{tikzcd}}}
La condition de fonctorialité en $X$ \inserted{et en $F$} \deleted{\unsure{nous ramène au}} permet de nous localiser dans $F$ et sur $X$, ce qui nous ramène au cas où $F\to X$ est une imm. ouverte, puis un isom, qui est p.ticulier de la $2^{\text{ième}}$ condition énoncée dans 1).

\figure{\begin{tikzpicture}[x=1cm,y=1cm]\draw (0,1.5)..controls(1.5,1.9)and(2,1.5)..(1,1)..controls(-.3,.15)and(-.3,-.1)..(1.45,0);\end{tikzpicture}}
Supposons $\operatorname{Fr}_X=\mathrm{id}$, (ce qui signifie que $X$ est défini par un espace loc. \unsure{quasi tt} discontinu -- pas néc. séparé -- muni du faisceau des fonctions continues à valeurs dans $\mathbf F_p$ $\cdots$) p.ex. $X=\operatorname{Spec}\mathbf F_p$. Alors on a un isom encore plus trivial $\operatorname{Fr}_X^*(F)=\mathrm{id}_X^*(F)\simeq F$, et $\operatorname{Fr}_{F/X}$ devient un \uline{autom} de $F$. Ce dernier n'est pas \deleted{\unsure{du}} \unsure{tjrs} l'identité. Il est l'identité \unsure{ssi} $F$ est constant sur $X=\operatorname{Spec}\mathbf F_p$
\mspage{31}
